% Alzheimer's Disease — see issue #26.
\subsection{Alzheimer's Disease}
\label{subsec:alzheimer}

Alzheimer's disease (AD) is a progressive neurodegenerative disorder that begins with memory impairment and ends in dementia. Treatment is effective only in the early stage, where it can control the disease but not cure it \parencite{nibila2026artificial}.

The disease begins at the cellular level. Amyloid-beta accumulates as plaques outside neurons and tau as tangles inside them, followed by neurodegeneration, the loss of synapses and neurons \parencite{jack2018niaaa}. This loss is not spread evenly: tangle pathology appears first in the entorhinal region, then reaches the hippocampus, and only later spreads across the neocortex \parencite{braak1991neuropathological}.

Clinically, the disease is described as a continuum that is usually divided into three groups \parencite{jack2018niaaa}, summarized in \cref{tab:ad_stages}.

% Clinical groups used in AD classification — used in sections/theoretical_framework/alzheimer.tex.
\begin{table}[t]
  \centering
  \caption{Cognitive stages used as labels in Alzheimer's disease classification.}
  \label{tab:ad_stages}
  \small
  \begin{tabular}{@{}p{0.12\columnwidth}p{0.73\columnwidth}@{}}
    \toprule
    Label & Description \\
    \midrule
    CN  & Cognitively normal (unimpaired). \\
    MCI & Mild cognitive impairment: measurable decline that does not yet prevent daily life; some, but not all, people progress to dementia. \\
    AD  & Dementia due to Alzheimer's disease: decline severe enough to interfere with daily life, attributed to Alzheimer's disease. \\
    \bottomrule
  \end{tabular}
\end{table}


Diagnosis combines clinical and neuropsychological assessment with biomarkers, measurable indicators of the disease:
\begin{itemize}
  \item \textbf{Cerebrospinal fluid (CSF):} amyloid-beta 42 and total tau are established markers \parencite{nibila2026artificial}.
  \item \textbf{Blood tests:} blood-based amyloid tests are less invasive but less accurate than CSF \parencite{nibila2026artificial}.
  \item \textbf{Genetic testing:} gives information about the risk of developing the disease \parencite{nibila2026artificial}.
  \item \textbf{Structural MRI:} shows the consequence of neurodegeneration as atrophy, with hippocampal atrophy the best-known sign \parencite{frisoni2010clinical}.
  \item \textbf{PET:} measures function rather than shape; FDG-PET shows reduced glucose metabolism and amyloid PET shows plaque deposition, but its cost, limited availability and radiation limit its use \parencite{nibila2026artificial}.
\end{itemize}

Atrophy follows the pattern of the pathology. It affects first the entorhinal cortex and the hippocampus, then the rest of the medial temporal lobe, including the amygdala, and later the lateral temporal and parietal cortex and the posterior cingulate; enlarged ventricles reflect the general loss of tissue \parencite{braak1991neuropathological,frisoni2010clinical}.
